\documentclass[a4paper,11pt,AutoFakeBold]{ctexart}
\usepackage{array}
\usepackage{amsmath}
\usepackage{amssymb}
\usepackage{array}
% 关于中文文献引用的格式参数设置请参考
% https://github.com/hushidong/biblatex-gb7714-2015
\usepackage[backend=biber,style=gb7714-2015]{biblatex}
\usepackage{fancyhdr}
\usepackage[margin=1in]{geometry}
\usepackage{graphicx}
\usepackage[hidelinks]{hyperref}
\usepackage{listings}
\usepackage{minted}
\usepackage{tabularx}
\usepackage{url}
\usepackage[dvipsnames]{xcolor}

\usepackage{amsmath,amssymb}
\usepackage{tikz}
\usepackage{tikz-3dplot}
\usetikzlibrary{calc,arrows.meta,angles,quotes,decorations.pathreplacing,positioning}


% 设置页眉，从第二页开始
\pagestyle{fancy}
\fancyhead[L]{王锦添}
\fancyhead[C]{缪子反常磁矩实验 Toy MC 模拟与数据拟合}
\fancyhead[R]{2026年4月}
\fancyfoot[C]{\thepage}
\renewcommand{\headrulewidth}{1pt}

% 定义行距=1.25倍
\linespread{1.25}

% 定义英文字体
\setmainfont{Times New Roman}
% 定义中文字体
\setCJKmainfont{FandolSong}
% 定义生僻字处理，当文字无法显示时前缀指令`\fallback`
\setCJKfamilyfont{Babel}{Microsoft YaHei} 
% or {SimHei}, or {KaiTi}
\newcommand{\fallback}{\CJKfamily{Babel}}

% 定义自定义命令
\newcommand{\spin}{\text{spin}}

% 设置一级标题左对齐
\ctexset{
  section={
    format+ =\raggedright
  }
}

% 定义常见软链颜色
\hypersetup{
  colorlinks = true,
  urlcolor = CadetBlue,
  linkcolor = Cerulean,
  citecolor = Maroon
}

% 定义行内代码格式
\definecolor{light-gray}{rgb}{0.96, 0.96, 0.96}
\NewDocumentCommand{
  \codeword}{v}{%
    \colorbox{light-gray}{
      \texttt{\textcolor{Black}{#1}
    }
  }%
}

% 定义引用源
\addbibresource{ref.bib}

\title{\textbf{缪子反常磁矩实验 Toy MC 模拟与数据拟合}}
\author{王锦添}
\date{2026年4月}












\begin{document}

\maketitle
%\thispagestyle{empty}   % 让标题页不显示页眉页脚

\pagestyle{fancy}       % 从下一页开始再启用页眉页脚

\section{引言}

标准模型（Standard Model）是现代物理学迄今为止最成功的理论之一\cite{BOSTON2023}，然而它并非完美无缺：在预测缪子的反常磁矩方面，标准模型的理论计算与实验测量之间存在着显著的差异，这被称为缪子 $g-2$ 问题。在美国费米实验室，缪子反常磁矩的精确测量实验已经进行多年，最新的结果显示了与标准模型预测之间的约4.2倍标准差的差距\cite{Li2022Muon}。因此，该实验的结果引发了物理学界的广泛关注和讨论，分析该实验数据的拟合过程对于理解这一现象至关重要。

这个程序旨在模拟缪子 $g-2$ 实验中的数据生成和分析过程，从模拟数据中提取物理参数，并通过拟合来验证理论模型。

\section{原理与实现}

实验中，极化的缪子被注入一个均匀的磁场中，其圆周运动的频率记作 $\omega_c$，而缪子自旋的进动频率记作 $\omega_s$。缪子的反常磁矩会导致 $\omega_s$ 与$\omega_c$ 之间存在一个差值，计算可得\cite{PhysRevA103}
\begin{equation}
\omega_a = \omega_s - \omega_c= - a_\mu \frac{q}{m_\mu} B
\end{equation}

其中 $a_\mu$ 是缪子的反常磁矩，定义为\cite{BOSTON2023} $a_\mu = \frac{g-2}{2}$。$q$ 是缪子的电荷，$m_\mu$ 是缪子的质量，$B$ 是磁场强度。通过测量 $\omega_a$ 与 $B$，我们可以间接地测量 $a_\mu$，从而检验标准模型的预测。



\subsection{正缪子衰变与正电子能量分布}

正缪子衰变为一个正电子和两个中微子 $\mu^+ \longmapsto e^+ + \nu_\mu + \bar{\nu}_e$ ，正电子的能量分布由衰变动力学决定。对于极化缪子，正电子概率密度函数由 Michel Spectrum 描述\cite{Fetscher2019MuonDecay}：
\begin{equation}
\frac{d^2 \Gamma}{dx\, d\cos\alpha}
=
\frac{G_F^2 m_\mu^5}{192\pi^3}
\left[
3 - 2x \pm P_\mu \cos\alpha (2x - 1)
\right] \cdot x^2
\end{equation}

等式左侧为单位能量分数 $x = {E_e}/{E_{max}}$ 和极化角（电子射出方向与缪子自旋方向的夹角） $\alpha$ 的微分衰变率，$G_F$ 是费米常数，$m_\mu$ 是缪子质量，$P_\mu$ 是缪子极化度。式子中的 $\pm$ 取正，因为实验中使用的是正缪子。记 $\cos \alpha = \Omega$，简化后正电子的能量分布可以表示为：
\begin{equation}
f(x,\Omega) = x^2 [3 - 2x + P_\mu \Omega (2x - 1)]
\end{equation}

\subsection{正电子能量计算}

考察缪子衰变的时刻。记缪子参考系为 $\Gamma'$，在该参考系中缪子处于静止状态。以缪子前进方向（动量 $\mathbf{p}_\mu$ 方向）为 $z$ 轴建立球坐标系，正电子的动量方向 $\mathbf{p}_e$ 由极角 $\theta$ 和方位角 $\phi$ 描述。此时，记缪子的自旋方向 $\mathbf{S}$ 与 $z$ 轴的夹角为 $\psi$。该角度可以由缪子相位计算得到，在代码中通过 $\psi = \omega_a T + \phi_0$ 计算，其中 $T$ 是缪子寿命（即缪子产生到此刻的时间），$\phi_0$ 是初始相位。则
\begin{equation}
\Omega = \cos\alpha = \mathbf{S} \cdot \mathbf{p}_e = \sin\psi \sin\theta \cos\phi + \cos\psi \cos\theta
\end{equation}


电子在 $\Gamma'$ 中可以获得的最高能量 $E'_{max} = (m_\mu^2 + m_e^2) / (2 m_\mu)$。 此时电子能量为 $E'=x E'_{max}$，动量为 $p'=\sqrt{E'^2 - m_e^2}$。根据狭义相对论，电子在 $\Gamma$ 中的能量 $E$ 可以通过以下公式计算\cite{Chu2021Muon}：
\begin{equation}
E = \gamma (E' + \beta p' \cos\theta )
\end{equation}


在实验中，我们设定了一个能量阈值 $E_{th}$，只有当正电子的能量 $E_e$ 满足 $E_e > E_{th}$ 时，才会被计数为一个事件。程序中通过比较计算得到的 $E$ 与 $E_{th}$ 来筛选事件，并将满足条件的事件记录下来用于后续分析。

\subsection{实验数据}

在程序中，为了接近实验条件，我选取了如下参数：

\begin{table}[htbp]
\centering
\begin{tabular}{llll}
\hline
\hline
参数 & 符号 & 数值 & 单位 \\
\hline
\hline
电子质量 & $m_e$ & $0.511$ & $\mathrm{MeV\, c^{-2}}$ \\
缪子质量 & $m_\mu$ & $105.658$ & $\mathrm{MeV\, c^{-2}}$ \\
缪子寿命 & $\tau_\mu$ & $2.2$ & $\mathrm{\mu s}$ \\
洛伦兹因子 & $\gamma$ & $29.3$ & - \\
缪子极化度\cite{PhysRevAccelBeams.24.044002}  & $P_\mu$ & $1.0$ & - \\
能量阈值\cite{BOSTON2023} & $E_{th}$ & $1700$ & $\mathrm{MeV}$  \\
缪子初始相位 & $\phi_0$ & $0$ & $\mathrm{rad}$ \\
缪子回旋频率\cite{BOSTON2023} & $\omega_c$ & $42.11$ & $\mathrm{rad\, s^{-1}}$ \\
反常自旋进动频率\cite{PhysRevLett.126.141801} & $\omega_a$ & $1.439$ & $\mathrm{rad\, s^{-1}}$ \\
\hline
\hline
\end{tabular}
\caption{模拟中使用的实验参数}
\end{table}


\begin{figure}[htbp]
\centering
\tdplotsetmaincoords{68}{120}
\begin{tikzpicture}[tdplot_main_coords,scale=3,>=Stealth]

    % Axes
    \draw[->,thick] (0,0,0) -- (1.45,0,0) node[below left] {$x$};
    \draw[->,thick] (0,0,0) -- (0,1.45,0) node[below right] {$y$};
    \draw[->,thick] (0,0,0) -- (0,0,1.45) node[above] {$z$};

    % Origin
    \fill (0,0,0) circle (0.015);
    \node[below left] at (0,0,0) {muon};

    % Muon spin vector before decay
    \coordinate (S) at (0.55,0.20,0.95);
    \coordinate (Sp) at (0.55,0.20,0);

    \draw[->,very thick,red!80!black] (0,0,0) -- (S) node[above right] {$\vec S$};
    \draw[dashed,red!60!black] (0,0,0) -- (Sp);
    \draw[dashed,red!60!black] (Sp) -- (S);

    % angle psi between S and z
    \tdplotdrawarc[->,red!70!black,thick]{(0,0,0)}{0.35}{0}{22}{anchor=south west}{$\psi$}

    % Text label
    \node[align=left,fill=white,inner sep=1pt] at (0.55,1.15,0.05)
    {\small \textbf{Before decay}\\[-1pt]
     \small muon spin direction fixed\\[-1pt]
     \small \(z\): reference axis};

\end{tikzpicture}
\caption{衰变前：只画出缪子的自旋方向 \(\vec S\) 和参考坐标轴。}
\end{figure}

\begin{figure}[htbp]
\centering
\tdplotsetmaincoords{68}{120}
\begin{tikzpicture}[tdplot_main_coords,scale=3,>=Stealth]

    % Axes
    \draw[->,thick] (0,0,0) -- (1.45,0,0) node[below left] {$x$};
    \draw[->,thick] (0,0,0) -- (0,1.45,0) node[below right] {$y$};
    \draw[->,thick] (0,0,0) -- (0,0,1.45) node[above] {$z$};

    % Origin / decay point
    \fill (0,0,0) circle (0.015);
    \node[below left] at (0,0,0) {decay point};

    % Spin vector retained as reference
    \coordinate (S) at (0.55,0.20,0.95);
    \coordinate (Sp) at (0.55,0.20,0);
    \draw[->,very thick,red!80!black] (0,0,0) -- (S) node[above right] {$\vec S$};
    \draw[dashed,red!60!black] (0,0,0) -- (Sp);
    \draw[dashed,red!60!black] (Sp) -- (S);

    % Positron momentum vector after decay
    \coordinate (P) at (0.95,0.62,0.52);
    \coordinate (Pp) at (0.95,0.62,0);
    \draw[->,very thick,blue!80!black] (0,0,0) -- (P) node[above left] {$\vec P$};
    \draw[dashed,blue!60!black] (0,0,0) -- (Pp);
    \draw[dashed,blue!60!black] (Pp) -- (P);

    % theta and phi
    \tdplotdrawarc[->,blue!70!black,thick]{(0,0,0)}{0.48}{0}{33}{anchor=west}{$\theta$}
    \tdplotdrawarc[->,blue!70!black,thick]{(0,0,0)}{0.33}{0}{34}{tdplot_main_coords,anchor=north}{$\phi$}

    % alpha between S and P
    \tdplotdrawarc[->,green!60!black,thick]{(0,0,0)}{0.62}{13}{58}{anchor=west}{$\alpha$}

    % Text label
    \node[align=left,fill=white,inner sep=1pt] at (0.55,1.15,0.05)
    {\small \textbf{After decay}\\[-1pt]
     \small emitted positron momentum \(\vec P\)\\[-1pt]
     \small \(\alpha\): angle between \(\vec S\) and \(\vec P\)};

\end{tikzpicture}
\caption{衰变后：画出正电子动量 \(\vec P\) 以及它与自旋 \(\vec S\) 的夹角 \(\alpha\)。}
\end{figure}







\newpage
\printbibliography[heading=bibliography,title=参考文献]

\end{document}
